\documentclass[../../main.tex]{subfiles}% 注意这里的文件路径不能用 ./main.tex ,否则用latexmk编译子文件会报错
\graphicspath{{\subfix{./image/}}{\subfix{../../image/}}} % 指定图片引用路径,后续可以直接使用图片文件名
% 如果图片存放在上级目录,则使用 ../../image/ 作为备用路径

% 例如:
% \begin{figure}[H]
% \centering
% \includegraphics[width=0.3\textwidth]{图.png}
% \caption{}
% \label{figure:图}
% \end{figure}
% 注意:上述\label{}一定要放在\caption{}之后,否则引用图片序号会只会显示??.

\begin{document}

\section{Young 测度}\label{section:ch02-Young测度}

设 $(X,\mathscr A,\mu)$ 为测度空间，$(X',\mathscr A')$ 为可测空间。设 $f:X\to X'$ 为 $(\mathscr A,\mathscr A')$-可测。利用\eqref{align:ch02-2.2}，立即可见映射
\begin{align}
\mu_f:\mathscr A'\to[0,\infty],\qquad \mu_f(E')\triangleq \mu(f^{-1}(E'))\quad(E'\in\mathscr A') \label{align:ch02-2.27}
\end{align}
为测度。

\begin{definition}\label{definition:ch02-像测度}
由\eqref{align:ch02-2.27}定义的测度 $\mu_f$ 称为 $\mu$ 在 $f$ 下的像。
\end{definition}

设 $(X_1\times X_2,\mathscr A_1\times\mathscr A_2)$ 为两个可测空间 $(X_1,\mathscr A_1)$ 与 $(X_2,\mathscr A_2)$ 的乘积，$\mu:\mathscr A_1\times\mathscr A_2\to[0,\infty]$ 为测度。由\cref{lemma:ch02-投影可测}，像测度
\begin{align}
\mu_{p_1}:\mathscr A_1\to[0,\infty],\qquad \mu_{p_1}(E_1)\triangleq \mu(p_1^{-1}(E_1))=\mu(E_1\times X_2)\quad(E_1\in\mathscr A_1), \label{align:ch02-2.28a}\\
\mu_{p_2}:\mathscr A_2\to[0,\infty],\qquad \mu_{p_2}(E_2)\triangleq \mu(p_2^{-1}(E_2))=\mu(X_1\times E_2)\quad(E_2\in\mathscr A_2) \label{align:ch02-2.28b}
\end{align}
良定义。

\begin{definition}\label{definition:ch02-投影测度}
由\eqref{align:ch02-2.28a}定义的测度 $\mu_{p_i}$ 称为 $\mu$ 在 $X_i$ 上的投影（$i=1,2$）。
\end{definition}

\begin{definition}\label{definition:ch02-Young测度}
设 $X$ 与 $Y$ 为局部紧 Hausdorff 空间，$\mu\in\mathscr R^+_f(X)$。
\begin{enumerate}[(i)]
\item 任意满足 $\nu(E\times Y)=\mu(E)$（对所有 $E\in\mathscr B(X)$）的测度 $\nu\in\mathscr R^+_f(X\times Y)$ 称为 Young 测度。
\item 设 $f:X\to Y$ 为 $(\mathscr B(X),\mathscr B(Y))$-可测。满足
\begin{align}
\nu(E\times F)=\mu(E\cap f^{-1}(F))\quad\text{对所有 }E\in\mathscr B(X),\ F\in\mathscr B(Y) \label{align:ch02-2.29}
\end{align}
的 Young 测度 $\nu$ 称为与 $f$ 相伴的 Young 测度。
\end{enumerate}
\end{definition}

由定义有 $\nu_{p_X}=\mu$，其中 $p_X$ 为 $X\times Y$ 到 $X$ 的投影。记 Young 测度的集合为 $\mathcal Y^+(X,\mathscr B(X),\mu;Y)\equiv\mathcal Y^+(X;Y)$。

\begin{remark}\label{remark:ch02-Young有界}
易见 $\mathcal Y^+(X;Y)$ 为 $\mathscr R^+_f(X\times Y)$ 的有界子集。事实上，对任意 $\nu\in\mathcal Y^+(X;Y)$，有
\begin{align*}
\|\nu\|_{\mathscr R_f(X\times Y)}=\nu(X\times Y)=\mu(X).
\end{align*}
故 $\mathcal Y^+(X;Y)$ 含于球面 $\partial B(0,\|\mu\|_{\mathscr R_f(X)})\subseteq\mathscr R^+_f(X\times Y)$。
\end{remark}

\begin{remark}\label{remark:ch02-Young图支撑}\label{remark:ch02-Young投影}
由\eqref{align:ch02-2.29}，与 $f$ 相伴的 Young 测度集中于 $f$ 的图像上，即集合 $\{(x,f(x))\mid x\in X\}$ 上（见\cref{definition:ch01-集中测度}）。又易见若 $f_1,f_2:X\to Y$ 为 $(\mathscr B(X),\mathscr B(Y))$-可测，$\nu_1,\nu_2$ 为相伴 Young 测度，则 $\nu_1=\nu_2$ 当且仅当 $f_1=f_2$ 在 $X$ 上 $\mu$-a.e.。事实上，若 $\nu_1=\nu_2$，则
\begin{align*}
\mu(X)=\nu_2(\{(x,f_2(x))\mid x\in X\})
=\nu_1(\{(x,f_2(x))\mid x\in X\})
=\mu(\{x\in X\mid f_1(x)=f_2(x)\}),
\end{align*}
故 $f_1=f_2$ $\mu$-a.e.。反向蕴含由\eqref{align:ch02-2.29}立即得到。
\end{remark}

最后提及下列结果。

\begin{lemma}\label{lemma:ch02-像测度积分}
\label{lemma:ch02-像测度}
设 $(X,\mathscr A,\mu)$ 为测度空间，$(X',\mathscr A')$ 为可测空间。设 $f:X\to X'$ 为 $(\mathscr A,\mathscr A')$-可测，$\mu_f$ 为 $\mu$ 在 $f$ 下的像。则对任意 $g\in L^1(X',\mathscr A',\mu_f)$，有 $g\circ f\in L^1(X,\mathscr A,\mu)$，且
\begin{align}
\int_F g\,\mathrm{d}\mu_f=\int_{f^{-1}(F)}g\circ f\,\mathrm{d}\mu\quad\text{对任意 }F\in\mathscr A'. \label{align:ch02-2.30}
\end{align}
\end{lemma}

\begin{proof}
设 $s\in\mathscr S(X')$，即 $s=\sum_{i=1}^n c_i\chi_{E'_i}$，其中 $c_1,\cdots ,c_n\in\mathbb R$，$\{E'_1,\cdots ,E'_n\}\subseteq\mathscr A'$ 为 $X'$ 的划分。因 $s\circ f=\sum_{i=1}^n c_i\chi_{f^{-1}(E'_i)}$，当 $g=s$ 时由\eqref{align:ch02-2.27}立即得\eqref{align:ch02-2.30}。再由 $\mathscr S(X')$ 在 $L^1(X',\mathscr A',\mu_f)$ 中的稠密性（见\cref{remark:ch02-简单函数稠密}）即得结论。
\end{proof}




\end{document}